\section*{الفصل \ref{ch:b3:orogenic-wedges} --- أحزمة الصدع الزاحف وأسافين التكوّن الجبلي}

\begin{solution}{exo:b3:orogenic-wedges:1}
$\sin\ang{30} = 0.5$: $K = 2\times0.5/0.5 = 2$ و$\sigma'_1/\sigma'_3 = 1 +
K = 3$.
\end{solution}

\begin{solution}{exo:b3:orogenic-wedges:2}
$\phi = \arctan0.6 = \ang{31.0}$، $\sin\phi = 0.514$، $K = 2.12$. على
قاعدة أفقية، $\alpha = 0.3/3.12 = \qty{0.0962}{rad} = \ang{5.5}$. وعلى
قاعدة مائلة بزاوية $\ang{3} = \qty{0.0524}{rad}$، $\alpha + \beta = 0.352/3.12 =
\qty{0.113}{rad} = \ang{6.5}$، إذن $\alpha = \ang{3.5}$.
\end{solution}

\begin{solution}{exo:b3:orogenic-wedges:3}
العمق $= (T - 10)/30$: الأباتيت \babelsublr{(U--Th)/He} من 1.3 إلى \qty{2.3}{km}؛ مسار الانشطار في الأباتيت
من 2.7 إلى \qty{3.7}{km}؛ مسار الانشطار في الزركون $270/30 = \qty{9.0}{km}$. عند
\qty{5}{km}، كانت الصخرة عند $10 + 150 = \qty{160}{\celsius}$: الزركون
\babelsublr{(U--Th)/He}، الذي يُغلق عند 160 إلى \qty{200}{\celsius}، على عمق
من 5.0 إلى \qty{6.3}{km}.
\end{solution}

\begin{solution}{exo:b3:orogenic-wedges:4}
الطول الأفقي $2/\tan\ang{30} = \qty{3.46}{km}$؛ على طول الانحدار
$2/\sin\ang{30} = \qty{4.0}{km}$. ترتفع الطبقة بمقدار ارتفاع المنحدر،
\qty{2}{km}. يجب أن يقطع خط القطع المنحدر كاملًا، فيكون الانزياح
مقداره \qty{4.0}{km}.
\end{solution}

\begin{solution}{exo:b3:orogenic-wedges:5}
$\phi = \arctan1.03 = \ang{45.8}$, $K = 5.08$; $\beta = \qty{0.105}{rad}$.
من أجل $\lambda = 0.4$: $\alpha + \beta = (0.105 + 0.85\times0.6)/(1 +
5.08\times0.6) = 0.615/4.05 = \qty{0.152}{rad} = \ang{8.7}$, $\alpha =
\ang{2.7}$. من أجل $\lambda = 0.9$: $0.190/1.51 = \qty{0.126}{rad} = \ang{7.2}$،
$\alpha = \ang{1.2}$. يؤدي رفع ضغط المائع إلى خفض \omterm{def:b3:orogenic-wedges:critical-taper}{الاستدقاق الحرج}:
يصبح الأسفين فوق الحرج، فينزلق إلى الأمام على قاعدته دون
تزايد في السماكة، وتتقدّم الجبهة حتى يعيده التراكم والتعرية
إلى الاستدقاق الجديد.
\end{solution}

\begin{solution}{exo:b3:orogenic-wedges:6}
$S = A/r = 6.0/1.5 = \qty{4.0}{km}$. تُقرأ بوصفها طية انفصال، $S = A/h =
6.0/3.0 = \qty{2.0}{km}$: نصف القيمة الحقيقية. ويجب التمييز بين القراءتين
بالاستناد إلى شكل الطية ووجود منحدر.
\end{solution}

\begin{solution}{exo:b3:orogenic-wedges:7}
$K = 2.12$ و$1 + K(1 - 0.4) = 2.27$. تعطي المعادلة $\mu_b(1 -
\lambda_b) = (\alpha + \beta)\times2.27 - \beta$، مع $\beta =
\qty{0.0349}{rad}$. جورا: $0.0611\times2.27 - 0.0349 = 0.10$. الجنوب:
$0.218\times2.27 - 0.0349 = 0.46$. بيرلي عند الضغط الهيدروستاتيكي:
$0.85\times0.6 = 0.51$. القاعدة الجنوبية قريبة من قيمة بيرلي، أما
قاعدة جورا فهي أضعف بنحو خمس مرات تقريبًا، كما يُتوقّع من سطح انفصال في
الهاليت، وهو ذو احتكاك منخفض، تحت سوائل غير مفرطة الضغط.
\end{solution}

\begin{solution}{exo:b3:orogenic-wedges:8}
$\kappa = \qty{e-6}{m^2/s} = \qty{31.6}{km^2/Myr}$، و$v =
\qty{1}{km/Myr}$: $\mathrm{Pe} = 30/31.6 = 0.951$. يكون المعامل $0.951/(1 -
e^{-0.951}) = 1.55$ يرفع التدرّج الحراري من 30 إلى \qty{46.5}{\celsius/km}.
من أجل \qty{65}{\celsius}: $1 - e^{-\mathrm{Pe}\,z/L} = (55/900)(1 - e^{-0.951})
= 0.0375$، إذن $\mathrm{Pe}\,z/L = 0.0382$ و$z = \qty{1.21}{km}$، في مقابل
$55/30 = \qty{1.83}{km}$ دون تعرية.
\end{solution}

\begin{solution}{exo:b3:orogenic-wedges:9}
$\mathrm dh/\mathrm dt = \qty{0.8}{km}/\qty{1.2}{Myr} = \qty{0.67}{mm/yr}$.
مع $f = 0.5$، $v = 0.67\times(1 - 0.5) = \qty{0.33}{mm/yr}$.
\end{solution}

\begin{solution}{exo:b3:orogenic-wedges:10}
مع $\alpha = 0$ و$u = 1 - \lambda$: $\beta(1 + Ku) = \beta + \mu_b u$،
إذن $\beta Ku = \mu_b u$ و$\beta^* = \mu_b/K$، مهما كانت $u$. بالنسبة إلى تايوان،
$0.85/5.08 = \qty{0.167}{rad} = \ang{9.6}$. وعلى قاعدة أشدّ انحدارًا، يكون
ميل السطح الحرج سالبًا: أيّ أسفين ذي سطح أفقي أو مائل نحو المقدّمة
يكون فوق الحرج أصلًا، وينزلق على قاعدته دون تزايد في السماكة. لا
يستطيع حزام الصدع الزاحف أن ينمو بشكله المعتاد هناك.
\end{solution}

\begin{solution}{exo:b3:orogenic-wedges:11}
عندما $\Lambda \to \infty$ و$v \to 0$، فإنّ $m \to 0$، إذن $e^{-mz_c} \to 1$،
والمنحنى الحراري مستقيم، $G_0 = G(z_c) = G$: $f \to (G - \Gamma)/G = 1 -
\Gamma/G = 1 - 6.5/30 = 0.78$. ويكون الميل عندئذٍ $v/(1 - f) = vG/\Gamma =
4.6v$. من أجل $v = \qty{1}{mm/yr}$، $v/2\kappa = \qty{0.0158}{km^{-1}}$. عند
\qty{5}{km}: $m = 0.0158 + \sqrt{0.0158^2 + 1.257^2} = \qty{1.273}{km^{-1}}$،
$e^{-1.273\times1.21} = 0.215$، $f = 40.0\times0.215/44.7 = 0.19$. عند
\qty{20}{km}: $m = 0.0158 + 0.3146 = \qty{0.330}{km^{-1}}$، $e^{-0.399} =
0.671$, $f = 40.0\times0.671/44.7 = 0.60$.
\end{solution}

\begin{solution}{exo:b3:orogenic-wedges:12}
المقطع مثلثٌ قاعدته $W$ وارتفاعه $W\tan(\alpha + \beta)$:
$A = W^2\tan(\alpha + \beta)/2$. مع $\mathrm dA/\mathrm dt = F_a$ و$A =
0$ عند $t = 0$، $A = F_at$، وهذا يعطي $W$. مع $\tan\ang{6} = 0.105$:
$W = \sqrt{2\times40\times2/0.105} = \qty{39}{km}$ بعد \qty{2}{Myr}؛ $t =
60^2\times0.105/(2\times40) = \qty{4.7}{Myr}$ للوصول إلى \qty{60}{km}.
\end{solution}

\begin{solution}{pb:b3:orogenic-wedges:1}
\textbf{1.} $4.5 - 1.5 = \qty{3.0}{km}$ ارتفاعًا؛ $3.0/\tan\ang{25} =
\qty{6.43}{km}$ طولًا أفقيًا؛ $3.0/\sin\ang{25} = \qty{7.10}{km}$ على امتداد
الانحدار.

\textbf{2.} بمقدار ارتفاع المنحدر، \qty{3.0}{km}.

\textbf{3.} $S = A/r = 16.5/3.0 = \qty{5.5}{km}$.

\textbf{4.} القيمة نفسها، \qty{16.5}{km^2}: في \omterm{def:b3:orogenic-wedges:thrust-folds}{طية انحناء الصدع} تكون المساحة الزائدة
لكل طبقة فوق السطح المستوي العلوي هي $Sr$. وبوصفها طية انفصال، $16.5/4.0 =
\qty{4.1}{km}$: أي أقلّ بنحو الربع.

\textbf{5.} $4\times(6.0 - 2.5) = \qty{14.0}{km}$، أي ضعفين ونصفًا قدر
الانزلاق على المنحدر الأمامي، \qty{5.5}{km}.

\textbf{6.} الكتلة الزاحفة الأمامية، أي الأقرب إلى المقدّمة: يتشكّل كل منحدر جديد
في الجدار المرتكز أمام المنحدر السابق. وكلّ كتلة، إذ تُرفع على منحدرها،
ترفع الصدع الزاحف السقفي والغطاء الذي فوقها كما ترفع \omterm{def:b3:orogenic-wedges:thrust-folds}{طية انحناء الصدع} جدارها
المعلّق، فيُطوى السقف فوق الكومة المتراكمة.

\textbf{7.} $\alpha = \arctan(2.1/60) = \ang{2.00}$; $\beta =
\arctan(4.2/60) = \ang{4.00}$؛ الاستدقاق \ang{6.0}.

\textbf{8.} $2.4 + 8.2 = \qty{10.6}{km}$، أي أقلّ من نحو \qty{15}{km}: القاعدة
ما زالت احتكاكية، وتنطبق عليها نظرية كولوم.

\textbf{9.} $\phi = \ang{45.8}$، $\sin\phi = 0.717$، $K = 5.08$. في الحالة الجافة:
$\alpha + \beta = (0.0699 + 0.85)/6.08 = \qty{0.151}{rad} = \ang{8.7}$,
$\alpha = \ang{4.7}$.

\textbf{10.} $\lambda = 1000/2500 = 0.40$: $(0.0699 + 0.51)/4.05 =
\qty{0.143}{rad} = \ang{8.2}$، $\alpha = \ang{4.2}$. أمّا الميل المرصود \ang{2.0}
سيجعل الإسفين دون الحرج بقوّة، فيسمك بدلًا من أن يتقدّم:
يجب أن تكون السوائل ذات ضغط زائد.

\textbf{11.} بأخذ $T = \qty{0.1049}{rad}$ و$u = 1 - \lambda$، يكون $T(1 + Ku) =
\beta + \mu_bu$ فيُعطي $u = (\beta - T)/(TK - \mu_b) = (0.0699 -
0.1049)/(0.5327 - 0.85) = 0.110$: $\lambda = 0.89$.

\textbf{12.} $\sigma_v = 2500\times9.81\times10\,600 = \qty{260}{MPa}$; $P =
0.89\times260 = \qty{231}{MPa}$؛ والضغط الهيدروستاتيكي $1000\times9.81\times10\,600 =
\qty{104}{MPa}$. تحمل السوائل \qty{127}{MPa} أكثر ممّا يحمله عمود
من الماء.

\textbf{13.} $(0.0699 + 0.85\times0.3)/(1 + 5.08\times0.3) = 0.325/2.52 =
\qty{0.129}{rad} = \ang{7.4}$، $\alpha = \ang{3.4}$. أمّا الإسفين، عند \ang{2.0}،
فهو الآن دون الحرج: يتوقّف عن التقدّم، ويسمك \omterm{def:b3:orogenic-wedges:thrust-geometry}{بالصدع الزاحف خارج التسلسل}
خلف جبهته حتى يبلغ ميل سطحه \ang{3.4}.

\textbf{14.} $F_a = \qty{10}{km/Myr}\times\qty{4.0}{km} = \qty{40}{km^2/Myr}$,
أي \qty{4.0e4}{m^3/yr} لكلّ كيلومتر من الحزام.

\textbf{15.} $\bar e = 40/60 = \qty{0.67}{km/Myr} = \qty{0.67}{mm/yr}$.

\textbf{16.} $40\times100 = \qty{4000}{km^3/Myr} = \qty{4.0e6}{m^3/yr}$، أو
$\qty{1.0e10}{kg/yr}$ (10 ملايين طنّ في السنة).

\textbf{17.} $\mathrm dW/\mathrm dt = F_a/(W\tan(\alpha + \beta)) =
40/(60\times0.105) = \qty{6.3}{km/Myr}$.

\textbf{18.} لا: في الحالة المستقرّة يجب أن تزيل التعرية كامل
التدفّق الوارد، فيبقى المعدّل المتوسط \qty{0.67}{mm/yr}. أمّا الصخر المضاف
أسفل الجزء الخلفي فيُرفع هناك، فستكون التعرية والاستخراج أسرع ما يكون
فوق الجزء الخلفي، حيث ستوجد أصغر أعمار التبرّد.

\textbf{19.} $v = \qty{0.667}{km/Myr}$: $\mathrm{Pe} = 0.667\times30/31.6 =
0.634$; $G_0 = 30\times0.634/(1 - e^{-0.634}) = 30\times1.35 =
\qty{40.5}{\celsius/km}$.

\textbf{20.} $1 - e^{-\mathrm{Pe}\,z/L} = (55/900)(1 - e^{-0.634}) = 0.0287$,
فإنّ $\mathrm{Pe}\,z/L = 0.0291$ و$z_c = \qty{1.38}{km}$؛ $G(z_c) =
40.5\,e^{-0.0291} = \qty{39.3}{\celsius/km}$.

\textbf{21.} $1.38/0.667 = \qty{2.07}{Myr}$، عمر العيّنة عند
\qty{2.00}{km}: عند متوسط الارتفاع يكون انحراف المتساوي الحرارة صفرًا.

\textbf{22.} $v/2\kappa = 0.667/63.1 = \qty{0.0106}{km^{-1}}$, $2\pi/\Lambda =
\qty{0.314}{km^{-1}}$: $m = 0.0106 + 0.3143 = \qty{0.325}{km^{-1}}$;
$e^{-0.325\times1.38} = 0.639$; $f = (40.5 - 6.5)\times0.639/39.3 = 0.552$.

\textbf{23.} يبلغ متوسط الارتفاع \qty{2.00}{km}، ومتوسط العمر
\qty{2.068}{Myr}؛ $\sum\Delta t\,\Delta h = 0.823$ و$\sum\Delta t^2 =
0.552$، فيكون ميل الخطّ \qty{1.49}{km/Myr}: معدّل ظاهريّ قدره
\qty{1.49}{mm/yr}.

\textbf{24.} $v = 1.49\times(1 - 0.552) = \qty{0.667}{mm/yr}$.

\textbf{25.} تتّفق \omterm{def:b3:orogenic-wedges:thermochronology}{الميقاتية الحرارية} وموازنة الدفق: الإسفين
في الحالة المستقرّة، ولو أُخذ الميل الظاهر على ظاهره لكان قد بالغ في تقدير المعدّل
بأكثر من ضعفين. ويبلغ معدّل الاستخراج لإسفين كيسترل
هو \textbf{\qty[text-series-to-math=true]{0.67}{mm/yr}}.
\end{solution}
